\documentclass[10pt,a4paper]{article}
\usepackage{amsmath,amssymb,amsthm,graphicx}
% Minimal retained preamble for S3 SIGMA 2022 article cards taken from
% sources/sigma-2022-012-a-quantum-0-infty-law/pinned-original-source.tex
% (doi:10.3842/SIGMA.2022.012, SIGMA 18 (2022), CC BY 4.0).
%
% Omitted and why: the publisher documentclass{sigma} block (source
% lines 1-7), the tikz packages (source lines 10-12, used only by the
% section-4 picture, which is not delivered) and the frontmatter
% commands - the sigma class is unavailable in the toolchain and the
% assembler supplies the article class. stmaryrd is kept: the author's
% intt macro (source line 25) uses \llbracket. allowdisplaybreaks is
% the author's own (source line 44).
% The author macro block (source lines 14-32, including the
% DeclareMathOperator for span) and the author theorem declarations
% (source lines 34-41, including the stray group brace on line 41) are
% copied verbatim below.
\usepackage{amsmath,amssymb,amsthm}
\usepackage{stmaryrd}
\allowdisplaybreaks
\DeclareMathOperator{\psan}{span}

\newcommand{\hilH}{\mathcal H}
\newcommand{\norm}[1]{\lVert #1\rVert}
\newcommand{\aver}[1]{\langle #1 \rangle}
\newcommand{\nN}{n\in \mathbb{N}}
\newcommand{\mN}{m\in \mathbb{N}}
\newcommand{\lN}{l\in \mathbb{N}}
\newcommand{\kN}{k\in \mathbb{N}}
\newcommand{\jN}{j\in \mathbb{N}}
\newcommand{\iN}{i\in \mathbb{N}}
\newcommand{\intt}[1]{\llbracket #1 \rrbracket}
\newcommand{\perpF}{{\perp _F}}
\newcommand{\expec}[1]{\mathbb{E}\left(#1\right)}
\newcommand{\algE}{{\mathcal E}}
\newcommand{\salg}{$\sigma$-algebra}
\newcommand{\rv}{random variable}
\newcommand{\Id}{\mathrm{Id}}
\newcommand{\ind}[1]{\mathbf{1}_{#1}}

\newtheorem{Theorem}{Theorem}[section]
\newtheorem{Corollary}[Theorem]{Corollary}
\newtheorem{Lemma}[Theorem]{Lemma}
\newtheorem{Proposition}[Theorem]{Proposition}
 { \theoremstyle{definition}
\newtheorem{Definition}[Theorem]{Definition}
\newtheorem{Example}[Theorem]{Example}
\newtheorem{Remark}[Theorem]{Remark} }
\title{Quantum 0-infinity law: divergence of the series of ratios is equivalent to totality of the random orthonormal sequence (Theorem 2.1)}
\author{Bauer, Michel}
\date{doi:10.3842/SIGMA.2022.012; CC BY 4.0}
\begin{document}
\maketitle
\setcounter{section}{2}
\setcounter{Theorem}{0}
\noindent\textit{Editorial note.} All delivered mathematical text is the author's own, copied from the pinned source file sources/sigma-2022-012-a-quantum-0-infty-law/pinned-original-source.tex; no statement or proof step was rewritten, reordered or removed.\par\medskip
\noindent\textit{Editorial note.} The counted span is the author Theorem environment theo:maindense (source lines 129-135); its proof is split by the author into two named proof environments (source lines 311-319 and 373-395, opening words 'Proof the first assertion' and 'Proof the second assertion' retained verbatim) and both are delivered as proof spans in the author's order.\par\medskip
\noindent\textit{Editorial note.} The definition and probability-setup contexts (source lines 94-120 and 127) and the U/V/W setup (181-188) sit before the statement; the sibling Lemma 4.2 (243-252) and Corollary 4.5 (289-295) are theorem-like contexts placed after the statement so the card keeps 2.1 (both print locally renumbered). Context line 98 is excluded because it cross-references the conclusions section, outside the card; the measure context stops at line 107 for the same reason.\par\medskip
\noindent\textit{Editorial note.} No bibliography fragment: no cite commands inside the delivered spans or contexts.\par\medskip
\noindent\textit{Editorial note.} Printed number anchored to the published PDF (SIGMA 18 (2022) 012, read from the stream-extracted text of sigma22-012.pdf, census pdf\_anchor\_evidence): the theorem family prints 2.1, 3.1 and 4.1-4.8 with the Corollary at 4.5 sitting mid-sequence between Lemma 4.4 and Lemma 4.6, confirming the sectionwise shared master counter Theorem declared in the author's newtheorem block (source lines 34-41). The pinned tex has five sections plus appendix A and its structure matches the published PDF.\par\medskip
\noindent\textit{Editorial note.} Theorem-family numbers inside the card are locally renumbered; every cross-reference is a symbolic ref and resolves inside the card. The source displays are unnumbered, so no equation numbering surface exists.\par\medskip
\noindent\textit{Editorial scope.} Single environment card cut from the SIGMA article doi:10.3842/SIGMA.2022.012 (sigma-2022-012-a-quantum-0-infty-law); statement and proof are the author's own text at source lines 129-135 and 311-319 and 373-395.\par\medskip
\noindent\textit{Allowed standard prerequisites (editorial).}\begin{itemize}
\item Matryoshka and probability setup. Delivered in context (source lines 94-120): the increasing dimension sequences, the n-Matryoshka F, the m-Matryoshka E in F with E\_k in F\_k, the Grassmannian measure mu and the recursive probability measure P on the set M of Matryoshkas, with the filtered space (M, algE, (algE\_k), P).. Load-bearing: fixes every object the statement and both proofs quantify over.
\item U/V/W decomposition setup. Delivered in context (source lines 181-188): the orthogonal projections U\_k(v), W\_k(v) and V\_k(v) and the three-term decomposition of v and of the Hilbert space.. Load-bearing: W\_k(v) and its projections are the quantities whose large-k behavior both proofs track.
\item Conditional-expectation recursion (Corollary 4.5). Delivered in context (source lines 289-295): E(norm W\_k(v)\textasciicircum{}2) = E(norm W\_\{i-1\}(v)\textasciicircum{}2) times the product of (1 - r\_j), bounded by norm(v)\textasciicircum{}2 times the same product.. Load-bearing: proof 1 reads the divergence case off it (source line 318) and proof 2 builds the lower-bound vectors on it (source line 376).
\item A.s.-to-zero characterization (Lemma 4.2). Delivered in context (source lines 243-252): W\_k(v) converges almost surely to zero iff the expected squared norms converge to zero.. Load-bearing: converts the expectation statement into the almost-sure totality claim (source line 316).
\end{itemize}\medskip
Let $\underline{n}:=(n_k)_{\kN}$ be an increasing sequence $0=n_0 \leq n_1\leq n_2\leq\cdots $ of integers with \mbox{$\lim_{k\to \infty}n_k=\infty$}.
An $\underline{n}$-Matryoshka $\underline{F}:=(F_k)_{\kN}$ is an increasing sequence $\{0\}=F_0\subset F_1\subset F_2\subset\cdots$ of subspaces of $\hilH$ with $\dim F_k:=n_k$ for $\kN$.
An $\underline{n}$-Matryoshka $\underline{F}$ is called total (in~$\hilH$) if $\cup_{\kN} F_k$ is dense in~$\hilH$.


In the framework of numerical computation of low-lying eigenstates, $F_k$ would be the $k^{\text{th}}$ trial space, the space in which the first $k$ levels are to be found.

We also fix an increasing sequence $\underline{m}:=(m_k)_{\kN}$ of integers such that $m_k\leq n_k$ for $\kN$ but $\lim_{k\to \infty}m_k=\infty$. We set $(\delta m)_k:=m_{k+1}-m_k$.
An $\underline{m}$-Matryoshka in $\underline{F}$ is an $\underline{m}$-Matryoshka $\underline{E}:=(E_k)_{\kN}$ such that $E_k \subset F_k$ for $\kN$. In~the sequel, if $\underline{E}$ is an $\underline{m}$-Matryoshka in~$\underline{F}$, $E_k^\perpF$ denotes the orthogonal complement of~$E_k$ in~$F_k$ (not in~$\hilH$).

We denote by $M_{\underline{F},\underline{m}}$, or simply $M$ when there is no risk of confusion, the set of $\underline{m}$-Matryoshkas in $\underline{F}$.
We recall that if $F$ is a finite-dimensional Hilbert space\footnote{In this study, a finite-dimensional subspace of a separable infinite-dimensional reference Hilbert space.} of dimension $n$ and $m\leq n$ there is an unique unitary invariant probability measure $\mu_F^m$ on the set of $m$-dimensional subspaces of~$F$ (a Grassmannian which we denote by $\mathbb{G}(F,m)$).


We use the measures $\mu_F^m$ to endow the set $M$ of $\underline{m}$-Matryoshkas in $\underline{F}$ with a probability measure $P_{\underline{F},\underline{m}}$, or simply $P$, defined recursively as follows:
\begin{itemize}\itemsep=0pt
\item $E_0=F_0=\{0\}$.
\item $E_0,\dots,E_k$ being chosen, $E_{k+1}$ is chosen as a random (in a unitary invariant way) subspace of $F_{k+1}$ containing $E_k$. The condition of unitary invariance can be rephrased as follows: $E_{k+1}$ is the orthogonal direct sum of $E_k$ with an $(\delta m)_k$-dimensional subspace of $E_k^\perpF$ chosen with $\mu_{E_k^\perpF}^{(\delta m)_k}$.
\end{itemize}

The recursive, Markovian, nature of this definition ensures that the appropriate analog of Kolmogorov's consistency criterion holds automatically, so that it extends to a probability measure on $\underline{m}$-Matryoshkas in $\underline{F}$, i.e., to infinite sequences.

More precisely, letting $\algE_k$ denote the \salg\ generated by the random subspaces $E_0,\dots,E_k$ and $\sigma \big(\cup_{\kN} \algE_k\big)$ the \salg\ they generate, there is an unique $\sigma \big(\cup_{\kN} \algE_k\big)$-probability measure on $M$ which has the right finite-dimensional marginals. We complete this probability measure \big(with the \salg\ $\sigma \big(\cup_{\kN} \algE_k\big)$ suitably enlarged to a \salg\ we denote by $\algE$ to include null sets\big), leading to the complete filtered probability space $\big(M,\algE, (\algE_k)_{\kN},P\big)$.

One basic question is under which conditions $\underline{E}$ is total \big(recall this means that $\cup_{\kN}E_k$ is dense in $\hilH$\big) with $P$-probability $1$. Denoting by $\overline{E}$ the closure of $\cup_{\kN}E_k$ in $\hilH$, the question is whether $\overline{E}^\perp$ \big(which is also the orthogonal of $\cup_{\kN}E_k$\big) reduces to $\{0\}$ (with $P$-probability $1$). This question has a simple and nice answer.

Fix an $\underline{m}$-Matryoshka $\underline{E}$ and let $v\in \hilH$. For $\kN$ denote by $U_k(v)$ (or simply $U_k$ when there is no risk of confusion) the orthogonal projection of $v$ on $E_k$ (a finite-dimensional, hence closed, subspace of $\hilH$) and set $W_k=W_k(v):=v-U_k(v)$, the projection of $v$ on the orthogonal complement of $E_k$ in $\hilH$. Set $V_0(v):=0$ and for $\kN^*$ set $V_k=V_k(v):=U_k(v)-U_{k-1}(v)$, which by construction is the orthogonal projection of $v$ on the orthogonal complement of $E_{k-1}$ in $E_k$. To summarize, for $\kN^*$, $v\in \hilH$ splits as a sum of three mutually orthogonal vectors
\[
v=U_{k-1}(v)+V_k(v)+W_k(v)
\]
corresponding to the orthogonal sum decomposition
\[
 \hilH = E_{k-1} \oplus \big(E_{k-1}^\perp  \cap  E_k\big) \oplus E_k^\perp.
 \]

\begin{Theorem} \label{theo:maindense}
For $\kN^*$ define $r_k:=\frac{m_k-m_{k-1}}{n_{k}-m_{k-1}}$ if $n_{k}-m_{k-1}>0$ and $r_k:=1$ otherwise.

 If the series $\sum_k r_k$ is divergent the space $\overline{E}^\perp$ is zero-dimensional $($i.e., $\underline{E}$ is total$)$ with pro\-ba\-bility $1$.

 If the series $\sum_k r_k$ is convergent the space $\overline{E}^\perp$ is infinite-dimensional with probability $1$.
\end{Theorem}

\begin{Lemma} \label{lem:basicfacts1} Let $v\in \hilH$.
\begin{enumerate}\itemsep=0pt
\item[$(i)$] As functions on $(M,\algE)$ with values in $\hilH$, $U_k(v)$, $V_k(v)$ and $W_k(v)$ are $\algE_k$-measurable for each $\kN$.

\item[$(ii)$] The sequence $W_k(v)$ converges almost surely to $0$ if and only if
\[
\lim_{k\to \infty} \expec{\norm{W_k(v)}^2}=0.
\]
\end{enumerate}
\end{Lemma}

\begin{Corollary} \label{cor:prodineq} Let $\iN^*$ and $v\in F_i$.
For $k \geq i$,
\[
\expec{\norm{W_k(v)}^2}=\expec{\norm{W_{i-1}(v)}^2}\prod_{j=i}^{k} (1-r_j) \leq \norm{v}^2\prod_{j=i}^{k} (1-r_j),
\]
with equality if $v$ belongs to the orthogonal complement of $F_{i-1}$ in $F_i$.
\end{Corollary}

\begin{proof}[Proof the first assertion of Theorem \ref{theo:maindense}]
By assumption, $\cup_{\iN}F_i$ is dense in $\hilH$. So $\cup_{\kN}E_k$ is dense in $\hilH$ if and only if it is dense in $\cup_{\iN}F_i$. Equivalently, $\cup_{\kN}E_k$ is dense in~$\hilH$ if and only if any $v$ in some $F_i$ is in the closure of $\cup_{\kN}E_k$, i.e., if and only if the sequence $(W_k(v))_{k\geq i}$ converges to $0$ for any $v$ in some $F_i$.

Using the linearity of the map $v\mapsto W_k(v)$, to prove that $\lim_{k\to \infty} W_k(v)=0$ for every $v\in \cup_{\iN}F_i$, it suffices to check this convergence when $v$ ranges over a basis of the finite-dimensional space $F_i$ for each $i$. We fix an orthonormal basis $(v_n)_{n\geq 1}$ in $\hilH$ such that $F_i:=\psan (v_n)_{n\in \intt{1,n_i}}$. Then $W_k(v_n))$ is well-defined for large enough $k$, i.e., for $k$ such that $n_k\geq n$: $\cup_{\kN}E_k$ is dense in $\hilH$ if and only if $\lim_{k \to \infty} W_k(v_n)=0$ for every $\nN^*$.

By $(ii)$ in Lemma \ref{lem:basicfacts1} we know that $\lim_{k\to \infty} W_k(v)=0$ holds almost surely for a given $v$ in some~$F_i$ if and only if $\lim_{k\to \infty} \expec{\norm{W_k(v)}^2}=0$. By the above remark, it is enough to let $v$ range over the countable set $(v_n)_{n\geq 1}$.

By Corollary \ref{cor:prodineq}, for $\iN^*$, $n\in \intt{n_{i-1}+1,n_i}$ and $k\geq i$, $\expec{\norm{W_k(v_n)}^2}= \prod_{i}^{k} (1-r_j)$. Thus $\cup_{\kN}E_k$ is dense in $\hilH$ if and only if the infinite product ``diverges'' to~$0$, i.e., if $\lim_{k\to \infty} \prod_{i}^{k} (1-r_j)\allowbreak=0$ for every $\iN^*$. As $r_j\in [0,1]$ for each $\jN ^*$, this is equivalent to the divergence of the series $\sum_j r_j$, concluding the proof.
\end{proof}

\begin{proof}[Proof the second assertion of Theorem \ref{theo:maindense}]
Our assumption now is that $\sum_k r_k$ is convergent. Hence, for large enough $i$, $s_i:=\prod_{k\geq i} (1-r_k) >0$, and $s_i \nearrow 1$ when $i\to \infty$.

By assumption, $\lim_{i\to \infty} n_i=\infty$. Thus $\mathbb{S}:=\{\iN^*, n_i-n_{i-1} >0\}$ is infinite. For each $i\in \mathbb{S}$ we choose a non-zero vector $v_i\in F_{i-1}^{\perp} \cap F_i$. By Corollary \ref{cor:prodineq}, $\expec{\norm{W_k(v_i)}^2}=\norm{v_i}^2\prod_{j=i}^{k} (1-r_j)$ for $k\geq i$. Taking the large $k$ limit, $\expec{\norm{\Phi v_i}^2}=\norm{v_i}^2s_i$ for $i\in \mathbb{S}$.

Fix $\varepsilon >0$, $\kN$ and $v\in F_k^\perp$. Write
\[ \expec{\norm{\Phi v}^2}=\expec{\norm{\Phi v}^2\ind{A_{k}(\varepsilon)}}+ \expec{\norm{\Phi v}^2\ind{A_{k}(\varepsilon)^c}}.\]
On $A_{k}(\varepsilon)$, $\norm{\Phi v}^2\leq \varepsilon^2 \norm{v}^2$ so
\[ \expec{\norm{\Phi v}^2\ind{A_{k}(\varepsilon)}}\leq \varepsilon^2 \norm{v}^2 P(A_{k}(\varepsilon)).\]
As $\norm{\Phi v}^2 \leq \norm{v}^2$ ($\Phi$ is an orthogonal projector!)
\[ \expec{\norm{\Phi v}^2\ind{A_{k}(\varepsilon)^c}}\leq \norm{v}^2P\big(A_{k}(\varepsilon)^c\big).\]

Thus we have proven that, for $\varepsilon >0$, $\kN$ and $v\in F_k^\perp$,
\[ \expec{\norm{\Phi v}^2} \leq \varepsilon^2 \norm{v}^2 P(A_{k}(\varepsilon)) +\norm{v}^2(1-P(A_{k}(\varepsilon)).\]
For $i\in \mathbb{S}$, $i>k$, it holds that $v_i\in F_{i-1}^{\perp} \cap F_i \subset F_k^\perp$, so putting things together
\[ \norm{v_i}^2s_i=\expec{\norm{\Phi v_i}^2} \leq \big(\varepsilon^2-1\big) \norm{v_i}^2 P(A_{k}(\varepsilon)) +\norm{v_i}^2.\]
Dividing by $ \norm{v_i}^2$ (which is $>0$ by construction) and rearranging yields
\[\big(1-\varepsilon^2\big)P(A_{k}(\varepsilon))\leq 1-s_i,\]
and letting $i\to \infty$ along $\mathbb{S}$ we obtain that $P(A_{k}(\varepsilon))=0$ for every $\varepsilon <1$, so $\cup_{\kN} A_{k}(\varepsilon)=\lim_{k\to \infty} A_{k}(\varepsilon)$ is a null set for every $\varepsilon <1$.

Thus $A:=\big\{\underline{E} \in M_{\underline{F},\underline{m}} \text{ such that } G_{\underline{E}}\text{ is finite-dimensional}\big\}$, a subset of $\cup_{\kN} A_{k}(\varepsilon)$ is also a null set. To summarize, if
$\sum_k r_k$ is convergent, the orthogonal $G_{\underline{E}}$ of $\cup_{\kN}E_k$ is infinite-dimensional with probability $1$, concluding the proof.
\end{proof}

\end{document}
